% ═══════════════════════════════════════════════════════════════
% 《第五维度》阅读导航图
% 编译方式：xelatex reading_map.tex
% ═══════════════════════════════════════════════════════════════
\documentclass{article}
\usepackage[paperwidth=26.5cm, paperheight=16.5cm, margin=8mm]{geometry}
\usepackage{xeCJK}
\setCJKmainfont{PingFang SC}
\usepackage{tikz}
\usetikzlibrary{calc, positioning, backgrounds, arrows.meta}
\pagestyle{empty}

% ── 配色（沿用全书 C+K 双通道体系） ──
\definecolor{cBlue}{cmyk}{1.0, 0, 0, 0.50}
\definecolor{cDC}{cmyk}{1.0, 0, 0, 0.30}
\definecolor{cAcc}{cmyk}{0.65, 0, 0, 0.05}
\definecolor{cFill}{cmyk}{0.12, 0, 0, 0}
\definecolor{cPale}{cmyk}{0.06, 0, 0, 0}
\definecolor{cG}{cmyk}{0, 0, 0, 0.58}
\definecolor{cLG}{cmyk}{0, 0, 0, 0.30}
\definecolor{cVLG}{cmyk}{0, 0, 0, 0.10}
% 篇章底色
\definecolor{bg1}{cmyk}{0.05, 0, 0, 0}
\definecolor{bg2}{cmyk}{0.08, 0, 0, 0}
\definecolor{bg3}{cmyk}{0.13, 0, 0, 0}
\definecolor{bg4}{cmyk}{0.09, 0, 0, 0}
\definecolor{bgF}{cmyk}{0.04, 0, 0, 0.01}

\begin{document}
\begin{tikzpicture}[x=1cm, y=1cm,
    chname/.style  = {font=\footnotesize,      text=cG,  align=center},
    chkey/.style   = {font=\tiny,              text=cLG, align=center},
    chq/.style     = {font=\tiny,              text=cAcc, align=center, text width=1.8cm},
    parttitle/.style = {font=\small\bfseries,  text=cBlue},
    partpoem/.style  = {font=\footnotesize,    text=cG},
    partdesc/.style  = {font=\tiny,            text=cLG},
]

% ═══════════════════════════════════════════════════
% 1. 山脊坐标  （x = 章节位置，y = 思想密度）
% ═══════════════════════════════════════════════════
\coordinate (r0)  at ( 0.0, 3.0);   % 序章
\coordinate (r1)  at ( 1.9, 2.5);   % 咫尺天涯
\coordinate (r2)  at ( 3.8, 3.3);   % 影分身术
\coordinate (r3)  at ( 5.7, 3.0);   % 数据的显微镜
\coordinate (r4)  at ( 7.6, 2.6);   % 遁甲天书
\coordinate (r5)  at ( 9.5, 4.3);   % 看不见的龙
\coordinate (r6)  at (11.4, 3.8);   % 月光宝盒
\coordinate (r7)  at (13.3, 5.8);   % 智能的阶梯  ★ 主峰
\coordinate (r8)  at (15.2, 3.3);   % 吸星大法
\coordinate (r9)  at (17.1, 5.0);   % 易筋经      ★ 第二峰
\coordinate (r10) at (19.0, 3.8);   % 丰裕社会
\coordinate (r11) at (20.9, 4.3);   % 变与不变
% 曲线首尾锚点
\coordinate (rL)  at (-1.0, 2.2);
\coordinate (rR)  at (21.7, 3.6);

% ═══════════════════════════════════════════════════
% 2. 篇章色带（背景层）
% ═══════════════════════════════════════════════════
\begin{scope}[on background layer]
  \fill[bg1]  ( 0.95, -3.0) rectangle ( 4.75, 7.6);   % 第一篇
  \fill[bg2]  ( 4.75, -3.0) rectangle ( 8.55, 7.6);   % 第二篇
  \fill[bg3]  ( 8.55, -3.0) rectangle (14.25, 7.6);   % 第三篇
  \fill[bg4]  (14.25, -3.0) rectangle (19.95, 7.6);   % 第四篇
  \fill[bgF]  (19.95, -3.0) rectangle (21.70, 7.6);   % 终章
\end{scope}
\foreach \xd in {0.95, 4.75, 8.55, 14.25, 19.95} {
  \draw[cVLG, line width=0.3pt] (\xd, -2.8) -- (\xd, 7.4);
}

% ═══════════════════════════════════════════════════
% 3. 山体填充 + 山脊线
% ═══════════════════════════════════════════════════
\fill[cFill, opacity=0.55]
  (-1.0, 0)
  -- plot[smooth, tension=0.50] coordinates
       {(rL)(r0)(r1)(r2)(r3)(r4)(r5)(r6)(r7)(r8)(r9)(r10)(r11)(rR)}
  -- (21.7, 0) -- cycle;

\draw[cBlue, line width=1.6pt]
  plot[smooth, tension=0.50] coordinates
    {(rL)(r0)(r1)(r2)(r3)(r4)(r5)(r6)(r7)(r8)(r9)(r10)(r11)(rR)};

% ═══════════════════════════════════════════════════
% 4. 章节节点
% ═══════════════════════════════════════════════════
\foreach \i/\lab in
  {0/序,1/一,2/二,3/三,4/四,5/五,6/六,7/七,8/八,9/九,10/十,11/终}
{
  \fill[white]              (r\i) circle (4.2mm);
  \draw[cBlue, line width=0.55pt] (r\i) circle (4.2mm);
  \node[font=\scriptsize\bfseries, text=cBlue] at (r\i) {\lab};
}
% 高亮两座山峰
\draw[cDC,  line width=1.3pt] (r7) circle (4.2mm);
\draw[cAcc, line width=1.1pt] (r9) circle (4.2mm);

% ═══════════════════════════════════════════════════
% 5. 山峰旗帜
% ═══════════════════════════════════════════════════
% ── 主峰：智能的阶梯 ──
\draw[cDC, line width=0.7pt] (13.3, 6.22) -- (13.3, 7.0);
\fill[cDC] (13.3, 7.0) -- (14.05, 6.75) -- (13.3, 6.50) -- cycle;
\node[font=\tiny\bfseries, text=cDC, anchor=west] at (14.1, 6.75)
  {智识之巅};

% ── 第二峰：易筋经 ──
\draw[cAcc, line width=0.7pt] (17.1, 5.42) -- (17.1, 6.1);
\fill[cAcc] (17.1, 6.1) -- (17.85, 5.85) -- (17.1, 5.60) -- cycle;
\node[font=\tiny\bfseries, text=cAcc, anchor=west] at (17.9, 5.85)
  {行动之巅};

% ═══════════════════════════════════════════════════
% 6. 章节名 + 关键词（挂在山脊下方）
% ═══════════════════════════════════════════════════

% ── 序章 ──
\node[chname, below=6mm of r0]  (n0)  {魔法世界};
\node[chkey,  below=0.5mm of n0] (k0)  {六道魔法 ・ 三条法则};
\node[chq,    below=0.5mm of k0]       {我们正走进\\什么样的世界？};

% ── 第一篇 ──
\node[chname, below=6mm of r1]  (n1)  {咫尺天涯};
\node[chkey,  below=0.5mm of n1] (k1)  {连接 ・ 延迟 ・ 共识};
\node[chq,    below=0.5mm of k1]       {速度的极限\\和代价是什么？};

\node[chname, below=6mm of r2]  (n2)  {影分身术};
\node[chkey,  below=0.5mm of n2] (k2)  {完备率\enspace 深\,>\,大};
\node[chq,    below=0.5mm of k2]       {数据要多\\还是要深？};

% ── 第二篇 ──
\node[chname, below=6mm of r3]  (n3)  {显微镜};
\node[chkey,  below=0.5mm of n3] (k3)  {迷彩服 ・ 定义权};
\node[chq,    below=0.5mm of k3]       {数据骗术\\长什么样？};

\node[chname, below=6mm of r4]  (n4)  {遁甲天书};
\node[chkey,  below=0.5mm of n4] (k4)  {可用不可见};
\node[chq,    below=0.5mm of k4]       {隐私如何与\\流动共存？};

% ── 第三篇 ──
\node[chname, below=6mm of r5]  (n5)  {看不见的龙};
\node[chkey,  below=0.5mm of n5] (k5)  {尺子召唤龙};
\node[chq,    below=0.5mm of k5]       {什么决定了\\AI进化的方向？};

\node[chname, below=6mm of r6]  (n6)  {月光宝盒};
\node[chkey,  below=0.5mm of n6] (k6)  {压缩即智能};
\node[chq,    below=0.5mm of k6]       {AI怎么学会\\"思考"？};

\node[chname, below=6mm of r7]  (n7)  {智能的阶梯};
\node[chkey,  below=0.5mm of n7] (k7)  {飞机与鸟 ・ 记忆 ・ 因果};
\node[chq,    below=0.5mm of k7]       {AI离真正智能\\还差什么？};

% ── 第四篇 ──
\node[chname, below=6mm of r8]  (n8)  {吸星大法};
\node[chkey,  below=0.5mm of n8] (k8)  {调用\,>\,拥有};
\node[chq,    below=0.5mm of k8]       {能力可调用后\\努力还有意义吗？};

\node[chname, below=6mm of r9]  (n9)  {易筋经};
\node[chkey,  below=0.5mm of n9] (k9)  {上下文 ・ 品味 ・ 目标};
\node[chq,    below=0.5mm of k9]       {如何训练自己的\\操作系统？};

\node[chname, below=6mm of r10] (n10) {丰裕社会};
\node[chkey,  below=0.5mm of n10](k10) {分布之外 ・ 无用之用};
\node[chq,    below=0.5mm of k10]      {智力丰裕后\\什么才稀缺？};

% ── 终章 ──
\node[chname, below=6mm of r11] (n11) {变与不变};
\node[chkey,  below=0.5mm of n11](k11) {简单 ・ 流动 ・ 进化};
\node[chq,    below=0.5mm of k11]      {技术剧变中\\什么值得坚守？};

% ═══════════════════════════════════════════════════
% 7. 概念桥梁（底层虚线箭头）
% ═══════════════════════════════════════════════════
% 完备率 → 自我完备率
\draw[-{Stealth[length=2mm, width=1.5mm]},
      cAcc, line width=0.45pt, densely dashed]
  (3.8, -0.1) .. controls (8, -1.3) and (13, -1.3) .. (17.1, -0.1);
\node[font=\tiny, text=cLG, fill=white, inner sep=1pt]
  at (10.5, -1.35) {完备率\;→\;自我完备率};

% 训练模型 → 训练自己
\draw[-{Stealth[length=2mm, width=1.5mm]},
      cAcc, line width=0.45pt, densely dashed]
  (11.4, -0.1) .. controls (13, -0.75) and (15.5, -0.75) .. (17.1, -0.1);
\node[font=\tiny, text=cLG, fill=white, inner sep=1pt]
  at (14.25, -0.85) {训练模型\;→\;训练自己};

% 序章悬念 → 第七章解答
\draw[-{Stealth[length=2mm, width=1.5mm]},
      cVLG, line width=0.45pt, densely dotted]
  (0.0, -0.1) .. controls (3, -2.1) and (10, -2.1) .. (13.3, -0.1);
\node[font=\tiny, text=cLG, fill=white, inner sep=1pt]
  at (6.65, -2.15) {AI会统治人类吗？→\;沿现有架构，不会};

% ═══════════════════════════════════════════════════
% 8. 篇章标题 + 诗句 + 副标题
% ═══════════════════════════════════════════════════
\node[parttitle]              at ( 0.0,  7.2) {序章};
\node[parttitle]              at ( 2.85, 7.2) {第一篇};
\node[parttitle]              at ( 6.65, 7.2) {第二篇};
\node[parttitle]              at (11.4,  7.2) {第三篇};
\node[parttitle]              at (17.1,  7.2) {第四篇};
\node[parttitle]              at (20.9,  7.2) {终章};

\node[partpoem]               at ( 2.85, 6.7) {于远峦可达};
\node[partpoem]               at ( 6.65, 6.7) {于微零可察};
\node[partpoem]               at (11.4,  6.7) {于过往可溯};
\node[partpoem]               at (17.1,  6.7) {于将来可塑};

\node[partdesc]               at ( 2.85, 6.25) {突破物理世界的法则};
\node[partdesc]               at ( 6.65, 6.25) {全知之眼的代价};
\node[partdesc]               at (11.4,  6.25) {以史为镜的炼金术};
\node[partdesc]               at (17.1,  6.25) {创世的权柄};

% ═══════════════════════════════════════════════════
% 9. 阅读体验标注（底部）
% ═══════════════════════════════════════════════════
\node[font=\tiny, text=cLG] at ( 2.85, -2.8) {轻松有趣};
\node[font=\tiny, text=cLG] at ( 6.65, -2.8) {烧脑实用};
\node[font=\tiny, text=cLG] at (11.4,  -2.8) {技术深水区};
\node[font=\tiny, text=cLG] at (17.1,  -2.8) {最贴近生活};

% ═══════════════════════════════════════════════════
% 10. 标题 + 图例
% ═══════════════════════════════════════════════════
\node[font=\large\bfseries, text=cBlue, anchor=south west]
  at (-1.0, 7.8) {《第五维度》阅读导航图};
\node[font=\tiny, text=cLG, anchor=south west]
  at (-1.0, -3.0) {山脊高度 $\approx$ 思想密度};

% 阅读方向箭头
\draw[-{Stealth[length=2.5mm]}, cLG, line width=0.4pt]
  (-0.5, -3.3) -- (21.5, -3.3);
\node[font=\tiny, text=cLG, anchor=east] at (21.5, -3.55) {阅读方向};

\end{tikzpicture}
\end{document}
